\documentclass[a4paper,12pt]{article}

\usepackage[T2A]{fontenc}
\usepackage[utf8]{inputenc}
\usepackage[russian]{babel}
\usepackage{amsmath,amssymb,mathtools}

% Надёжный кириллический sans-serif для pdfLaTeX.
\usepackage{DejaVuSans}
\renewcommand{\familydefault}{\sfdefault}

\usepackage{icomma}
\usepackage[a4paper,margin=2.1cm]{geometry}
\usepackage{microtype}

\usepackage{tikz}
\usetikzlibrary{calc,arrows.meta,positioning,shapes.geometric}

\usepackage{tabularx,booktabs}
\usepackage{enumitem}
\usepackage{hyperref}
\usepackage{parskip}
\usepackage{fancyhdr}
\usepackage{setspace}
\usepackage{xcolor}

\definecolor{accent}{HTML}{008080}
\definecolor{darkaccent}{HTML}{004D4D}
\definecolor{textgray}{HTML}{333333}

\hypersetup{
  colorlinks=true,
  linkcolor=darkaccent,
  urlcolor=darkaccent,
  pdftitle={Чек-лист ОГЭ: дроби, числовая прямая и сравнение чисел},
  pdfauthor={Лёвин Артём Александрович}
}

\setlength{\headheight}{25pt}
\onehalfspacing
\setlength{\parindent}{0pt}

\setlist[itemize]{
  leftmargin=1.6em,
  itemsep=0.35em,
  topsep=0.35em
}

\setlist[enumerate]{
  leftmargin=1.9em,
  itemsep=0.45em,
  topsep=0.4em
}

\renewcommand{\arraystretch}{1.35}

\pagestyle{fancy}
\fancyhf{}
\fancyhead[L]{\small\textcolor{textgray}{ОГЭ по математике}}
\fancyhead[R]{\small\textcolor{textgray}{София Хоменко}}
\fancyfoot[C]{\small\textcolor{textgray}{\thepage}}

\renewcommand{\headrulewidth}{0.4pt}
\renewcommand{\headrule}{%
  \hbox to\headwidth{%
    \color{accent}\leaders\hrule height \headrulewidth\hfill
  }%
}

\newcommand{\topic}[1]{%
  \vspace{0.3em}
  {\Large\bfseries\textcolor{darkaccent}{#1}}\par
  \vspace{0.3em}
  \noindent\textcolor{accent}{\rule{\linewidth}{0.7pt}}\par
  \vspace{0.5em}
}

\newcommand{\subtopic}[1]{%
  \vspace{0.6em}
  {\large\bfseries\textcolor{darkaccent}{#1}}\par
  \vspace{0.15em}
}

\newcommand{\important}[1]{%
  \textbf{\textcolor{darkaccent}{#1}}%
}

\tikzset{
  flowstart/.style={
    font=\small,
    ellipse,
    draw=darkaccent,
    line width=0.9pt,
    align=center,
    minimum width=3.5cm,
    minimum height=1cm,
    inner sep=5pt
  },
  flowaction/.style={
    font=\small,
    rounded corners=4pt,
    rectangle,
    draw=accent,
    line width=0.9pt,
    align=center,
    text width=8.9cm,
    minimum height=1.05cm,
    inner sep=6pt
  },
  flowdecision/.style={
    font=\small,
    diamond,
    aspect=2.25,
    draw=darkaccent,
    line width=0.9pt,
    align=center,
    text width=4cm,
    inner sep=3pt
  },
  flowend/.style={
    font=\small,
    ellipse,
    draw=darkaccent,
    line width=0.9pt,
    align=center,
    minimum width=4.2cm,
    minimum height=1cm,
    inner sep=5pt
  },
  flowarrow/.style={
    -{Stealth[length=3mm,width=2mm]},
    line width=0.9pt,
    draw=textgray
  },
  flowlabel/.style={
    font=\small,
    fill=white,
    inner sep=1.5pt
  }
}

\begin{document}

% =========================================================
% ТИТУЛЬНЫЙ ЛИСТ
% =========================================================

\hypersetup{pageanchor=false}

\begin{titlepage}
\thispagestyle{empty}
\centering

\vspace*{2cm}

{\Large\textcolor{darkaccent}{\textbf{Чек-лист}}\par}

\vspace{0.8cm}

{\huge\bfseries
Дроби, числовая прямая\\[0.25em]
и сравнение чисел
\par}

\vspace{0.7cm}

{\large Подготовка к ОГЭ по математике\par}

\vfill

{\large\textbf{София Хоменко}\par}

\vspace{0.4cm}

{\normalsize
Лёвин Артём Александрович\\
эксклюзивно для Софии Хоменко
\par}

\vspace{0.7cm}

{\normalsize \today\par}

\vfill

\begin{tikzpicture}[remember picture,overlay]
  \draw[
    accent,
    line width=1.1pt
  ]
  ([xshift=1.8cm,yshift=1.5cm]current page.south west)
  .. controls
  ([xshift=6cm,yshift=2.05cm]current page.south west)
  and
  ([xshift=-6cm,yshift=0.95cm]current page.south east)
  ..
  ([xshift=-1.8cm,yshift=1.5cm]current page.south east);
\end{tikzpicture}

\end{titlepage}

\hypersetup{pageanchor=true}

% =========================================================
% ОСНОВНОЙ ЧЕК-ЛИСТ
% =========================================================

\topic{1. Обыкновенные и смешанные дроби}

\subtopic{1.1. Приведение к общему знаменателю}

Для сложения и вычитания дробей знаменатели должны быть одинаковыми:
\[
\frac{a}{b}\pm\frac{c}{d}
=
\frac{ad\pm bc}{bd},
\qquad
b\ne0,\quad d\ne0.
\]

Формула работает для любых ненулевых знаменателей. При этом знаменатель
\(bd\) иногда оказывается больше необходимого. Если НОК знаменателей
находится быстро, обычно выгодно использовать его.

\important{Пример.}
\[
\frac{4}{25}+\frac{15}{4}
=
\frac{16}{100}+\frac{375}{100}
=
\frac{391}{100}
=
3{,}91.
\]

Для знаменателей \(30\) и \(42\):
\[
30=2\cdot3\cdot5,
\qquad
42=2\cdot3\cdot7.
\]

Поэтому
\[
\text{НОК}(30,42)
=
2\cdot3\cdot5\cdot7
=
210.
\]

Дополнительные множители:
\[
30\cdot7=210,
\qquad
42\cdot5=210.
\]

\subtopic{1.2. Смешанная дробь сначала становится неправильной}

\[
m\frac{n}{q}
=
\frac{mq+n}{q},
\qquad
q\ne0.
\]

\important{Пример.}
\[
6\frac12-\frac{47}{10}
=
\frac{13}{2}-\frac{47}{10}
=
\frac{65}{10}-\frac{47}{10}
=
\frac{18}{10}
=
1{,}8.
\]

\subtopic{1.3. Деление дробей}

\[
\frac{a}{b}:\frac{c}{d}
=
\frac{a}{b}\cdot\frac{d}{c}
=
\frac{ad}{bc},
\]
где
\[
b\ne0,
\qquad
d\ne0,
\qquad
c\ne0.
\]

Перед умножением полезно сокращать общие множители числителя
и знаменателя.

\important{Пример.}
\[
\frac{12}{5}:\frac{15}{2}
=
\frac{12}{5}\cdot\frac{2}{15}
=
\frac{4}{5}\cdot\frac{2}{5}
=
\frac{8}{25}
=
0{,}32.
\]

Целое число можно считать дробью со знаменателем \(1\):
\[
2=\frac21.
\]

Поэтому
\[
\frac{\frac53}{2}
=
\frac{\frac53}{\frac21}
=
\frac56,
\]
а также
\[
\frac{5}{\frac32}
=
\frac{\frac51}{\frac32}
=
\frac{10}{3}.
\]

\subtopic{1.4. Четырёхэтажная дробь}

Если в числителе и знаменателе стоят дроби, получаем:
\[
\frac{\frac{a}{b}}{\frac{c}{d}}
=
\frac{ad}{bc},
\qquad
b\ne0,\quad c\ne0,\quad d\ne0.
\]

Удобная мнемоника:
\[
\frac{\frac{a}{b}}{\frac{c}{d}}
\quad\longrightarrow\quad
\frac{\text{произведение внешних}}
     {\text{произведение внутренних}}.
\]

В основе правила остаётся обычное деление дробей:
делитель заменяется обратной дробью.

% =========================================================
% БЛОК-СХЕМА 1
% =========================================================

\clearpage
\thispagestyle{plain}

\begin{center}
{\Large\bfseries
\textcolor{darkaccent}{Алгоритм: действия с дробями}}
\end{center}

\vspace{0.2cm}

\begin{center}
\begin{tikzpicture}[node distance=6mm and 12mm]

\node[flowstart] (start) {Начало};

\node[
  flowaction,
  below=of start
] (convert) {
  Смешанные числа перевести в неправильные дроби.\\
  Десятичные дроби при необходимости перевести в обыкновенные.
};

\node[
  flowdecision,
  below=8mm of convert
] (op) {
  Какое действие?
};

\node[
  flowaction,
  below left=9mm and 8mm of op,
  text width=5cm
] (add) {
  Сложение или вычитание:
  привести к общему знаменателю.
  Если НОК находится быстро --- использовать НОК;
  иначе допустима формула
  \(\frac{ad\pm bc}{bd}\).
};

\node[
  flowaction,
  below right=9mm and 8mm of op,
  text width=5cm
] (div) {
  Деление:
  проверить, что делитель не равен нулю;
  заменить деление умножением на обратную дробь.
};

\node[
  flowaction,
  below=12mm of op,
  yshift=-39mm
] (reduce) {
  Сократить общие множители там, где сокращение разрешено,
  затем выполнить вычисления.
};

\node[
  flowdecision,
  below=7mm of reduce
] (format) {
  Нужен десятичный ответ?
};

\node[
  flowaction,
  below left=8mm and 8mm of format,
  text width=5cm
] (decimal) {
  Перевести результат в десятичную дробь.
};

\node[
  flowaction,
  below right=8mm and 8mm of format,
  text width=5cm
] (exact) {
  Оставить сокращённую обыкновенную дробь.
};

\node[
  flowend,
  below=12mm of format,
  yshift=-27mm
] (end) {
  Ответ
};

\draw[flowarrow] (start) -- (convert);
\draw[flowarrow] (convert) -- (op);

\draw[flowarrow]
(op)
-| node[flowlabel,pos=0.22] {\(+\) / \(-\)}
(add);

\draw[flowarrow]
(op)
-| node[flowlabel,pos=0.22] {:}
(div);

\draw[flowarrow] (add) |- (reduce);
\draw[flowarrow] (div) |- (reduce);

\draw[flowarrow] (reduce) -- (format);

\draw[flowarrow]
(format)
-| node[flowlabel,pos=0.22] {да}
(decimal);

\draw[flowarrow]
(format)
-| node[flowlabel,pos=0.22] {нет}
(exact);

\draw[flowarrow] (decimal) |- (end);
\draw[flowarrow] (exact) |- (end);

\end{tikzpicture}
\end{center}

\clearpage

% =========================================================
% ДЕСЯТИЧНЫЕ ДРОБИ
% =========================================================

\topic{2. Десятичные дроби и деление}

\subtopic{2.1. Как убрать запятую из делителя}

В частном можно умножить делимое и делитель
на одно и то же число \(10^k\):
\[
\frac{x}{y}
=
\frac{x\cdot10^k}{y\cdot10^k},
\qquad
y\ne0.
\]

Выбирайте минимальное \(k\), при котором делитель становится целым.

\important{Пример.}
\[
1{,}25:0{,}002
=
\frac{1{,}25}{0{,}002}
=
\frac{1250}{2}
=
625.
\]

\important{Контрольный пример.}
\[
2{,}4:1{,}08
=
\frac{240}{108}
=
\frac{20}{9}
=
2{,}(2).
\]

\subtopic{2.2. Когда удобно доводить знаменатель до \(10\), \(100\), \(1000\)}

Если после сокращения знаменатель легко превратить
в степень десяти, вычисление существенно упрощается:
\[
\frac{8}{25}
=
\frac{32}{100}
=
0{,}32,
\]
\[
\frac{13}{20}
=
\frac{65}{100}
=
0{,}65.
\]

После сокращения обыкновенная дробь имеет конечную десятичную запись
тогда и только тогда, когда в разложении её знаменателя остаются
только простые множители \(2\) и \(5\).

\subtopic{2.3. Полезные дроби, которые стоит узнавать сразу}

\begin{table}[ht]
\centering
\caption{Часто встречающиеся дроби}

\begin{tabularx}{\linewidth}{
  @{}
  >{\centering\arraybackslash}X
  >{\centering\arraybackslash}X
  >{\centering\arraybackslash}X
  >{\centering\arraybackslash}X
  @{}
}
\toprule
\(\frac12=0{,}5\) &
\(\frac14=0{,}25\) &
\(\frac38=0{,}375\) &
\(\frac15=0{,}2\)
\\
\midrule
\(\frac13=0{,}(3)\) &
\(\frac16=0{,}1(6)\) &
\(\frac23=0{,}(6)\) &
\(\frac56=0{,}8(3)\)
\\
\bottomrule
\end{tabularx}
\end{table}

% =========================================================
% ЧИСЛОВАЯ ПРЯМАЯ
% =========================================================

\topic{3. Числовая прямая и подстановка}

\subtopic{3.1. Сначала оцените положение точки}

Для точки на числовой прямой сначала определите:

\begin{itemize}
  \item знак числа;
  \item между какими соседними числами оно расположено;
  \item ближе к какому из них находится;
  \item при двух точках --- какая дальше от нуля, то есть у какой больше модуль.
\end{itemize}

\begin{center}
\begin{tikzpicture}[x=1.55cm,y=1cm]

  \draw[
    -{Stealth[length=2.5mm]},
    line width=0.9pt
  ]
  (-3.4,0) -- (3.5,0);

  \foreach \x in {-3,-2,-1,0,1,2,3}{
    \draw
    (\x,0.09)
    --
    (\x,-0.09)
    node[below=4pt]{\small \(\x\)};
  }

  \fill[accent]
  (-2.4,0)
  circle (2.2pt)
  node[
    above=6pt,
    text=darkaccent
  ]{\small \(a\approx-2{,}4\)};

\end{tikzpicture}
\end{center}

\subtopic{3.2. Отрицательное число при подстановке берите в скобки}

Если
\[
a\approx-2{,}4,
\]
записывайте:
\[
-1-a
=
-1-(-2{,}4)
=
1{,}4.
\]

Скобки особенно важны при:

\begin{itemize}
  \item вычитании;
  \item умножении;
  \item возведении в степень.
\end{itemize}

Если варианты ответа содержат неравенства,
каждое утверждение проверяйте вычислением.

Если формат задания требует номер варианта,
в ответ записывается именно
\important{номер подходящего ответа}.

\subtopic{3.3. Расстояние от нуля и модуль}

Чем дальше точка расположена от \(0\),
тем больше модуль соответствующего числа:
\[
|a|>|b|
\quad\Longleftrightarrow\quad
\text{точка }a\text{ дальше от }0
\text{, чем точка }b.
\]

Знак и модуль характеризуют число по-разному:
\[
-5<2,
\qquad
|-5|>|2|.
\]

% =========================================================
% БЛОК-СХЕМА 2
% =========================================================

\clearpage
\thispagestyle{plain}

\begin{center}
{\Large\bfseries
\textcolor{darkaccent}{Алгоритм: утверждения по числовой прямой}}
\end{center}

\vspace{0.2cm}

\begin{center}
\begin{tikzpicture}[node distance=6mm]

\node[flowstart] (start2) {
  Начало
};

\node[
  flowaction,
  below=of start2
] (estimate) {
  Оценить каждую отмеченную точку:
  знак, соседние числа, приблизительное значение.
};

\node[
  flowaction,
  below=of estimate
] (distance) {
  Если точек несколько,
  сравнить их положение относительно \(0\) и модули.
};

\node[
  flowaction,
  below=of distance
] (substitute) {
  Подставить значения в очередное утверждение.
  Отрицательные числа записывать в скобках.
};

\node[
  flowaction,
  below=of substitute
] (calc) {
  Аккуратно вычислить левую часть утверждения.
};

\node[
  flowdecision,
  below=8mm of calc
] (truth) {
  Утверждение верно?
};

\node[
  flowaction,
  below=10mm of truth
] (mark) {
  Зафиксировать результат проверки
  и сопоставить его с формулировкой задания.
};

\node[
  flowdecision,
  below=8mm of mark
] (more) {
  Остались варианты?
};

\node[
  flowend,
  below=8mm of more
] (end2) {
  Записать требуемый номер ответа
};

\draw[flowarrow] (start2) -- (estimate);
\draw[flowarrow] (estimate) -- (distance);
\draw[flowarrow] (distance) -- (substitute);
\draw[flowarrow] (substitute) -- (calc);
\draw[flowarrow] (calc) -- (truth);

\draw[flowarrow]
(truth.west)
-- ++(-1.3cm,0)
node[
  flowlabel,
  above,
  pos=0.45
]{да}
|- (mark.west);

\draw[flowarrow]
(truth.east)
-- ++(1.3cm,0)
node[
  flowlabel,
  above,
  pos=0.45
]{нет}
|- (mark.east);

\draw[flowarrow] (mark) -- (more);

\draw[flowarrow]
(more)
--
node[flowlabel]{нет}
(end2);

\draw[flowarrow]
(more.east)
-- ++(4.5cm,0)
node[
  flowlabel,
  above,
  pos=0.18
]{да}
|- ([xshift=7mm]substitute.east)
-- (substitute.east);

\end{tikzpicture}
\end{center}

\clearpage

% =========================================================
% СРАВНЕНИЕ ЧИСЕЛ
% =========================================================

\topic{4. Сравнение чисел разных видов}

\subtopic{4.1. Универсальный приём: привести числа к удобной форме}

Для сравнения обыкновенных и десятичных дробей
часто удобно перевести всё в десятичную запись:
\[
\frac16=0{,}1(6),
\qquad
\frac14=0{,}25.
\]

Например,
\[
\frac16<0{,}2<\frac14.
\]

Бесконечную десятичную дробь обычно не требуется вычислять
с большой точностью. Достаточно такого количества цифр,
которое однозначно определяет порядок чисел.

\subtopic{4.2. Квадратные корни: зажать между соседними квадратами}

Для \(x\ge0\) найдите такие целые \(n\), что
\[
n^2<x<(n+1)^2.
\]

Тогда
\[
n<\sqrt{x}<n+1.
\]

\important{Пример 1.}
\[
49<53<64
\quad\Longrightarrow\quad
7<\sqrt{53}<8.
\]

\important{Пример 2.}
\[
64<73<81
\quad\Longrightarrow\quad
8<\sqrt{73}<9.
\]

Если требуется более точное сравнение,
для положительных величин можно сравнить квадраты.

Например:
\[
7{,}3^2=53{,}29>53.
\]

Следовательно,
\[
\sqrt{53}<7{,}3.
\]

\subtopic{4.3. Типичные ошибки}

\begin{itemize}

  \item
  Сокращать слагаемые через знак \(+\) или \(-\).
  Сокращение выполняется между множителями.

  \item
  При делении дробей забывать заменить делитель обратной дробью.

  \item
  Забывать ограничение
  \[
  \text{делитель}\ne0.
  \]

  \item
  При подстановке отрицательного числа терять скобки.

  \item
  Путать порядок чисел с порядком их модулей.

  \item
  Для квадратного корня выбирать числа,
  которые не являются квадратами соседних целых.
  Для \(\sqrt{53}\) правильные опорные значения:
  \[
  49=7^2,
  \qquad
  64=8^2.
  \]

  \item
  Вычислять десятичную запись с избыточной точностью,
  когда для сравнения достаточно нескольких первых цифр.

\end{itemize}

% =========================================================
% БЛОК-СХЕМА 3
% =========================================================

\clearpage
\thispagestyle{plain}

\begin{center}
{\Large\bfseries
\textcolor{darkaccent}{Алгоритм: сравнение дробей и корней}}
\end{center}

\vspace{0.2cm}

\begin{center}
\begin{tikzpicture}[node distance=6mm]

\node[flowstart] (s3) {
  Начало
};

\node[
  flowaction,
  below=of s3
] (take) {
  Рассматривать сравниваемые числа по одному.
};

\node[
  flowdecision,
  below=8mm of take
] (rootq) {
  Текущее число содержит квадратный корень?
};

\node[
  flowaction,
  below left=10mm and 5mm of rootq,
  text width=4.2cm
] (rootstep) {
  Если под корнем полный квадрат ---
  извлечь корень точно.
  Иначе зажать число между соседними квадратами
  и получить оценку.
};

\node[
  flowaction,
  below right=10mm and 5mm of rootq,
  text width=4.2cm
] (fracstep) {
  Обыкновенную дробь привести к удобному виду:
  чаще всего к десятичной записи.
  Десятичное число оставить как есть.
};

\node[
  flowaction,
  below=14mm of rootq,
  yshift=-43mm
] (record) {
  Записать сопоставимую числовую оценку
  с достаточной точностью.
};

\node[
  flowdecision,
  below=8mm of record
] (more3) {
  Остались числа?
};

\node[
  flowaction,
  below=8mm of more3
] (compare) {
  Сравнить полученные значения
  и расположить числа в требуемом порядке.
};

\node[
  flowend,
  below=of compare
] (e3) {
  Ответ
};

\draw[flowarrow] (s3) -- (take);
\draw[flowarrow] (take) -- (rootq);

\draw[flowarrow]
(rootq)
-| node[flowlabel,pos=0.22]{да}
(rootstep);

\draw[flowarrow]
(rootq)
-| node[flowlabel,pos=0.22]{нет}
(fracstep);

\draw[flowarrow] (rootstep) |- (record);
\draw[flowarrow] (fracstep) |- (record);

\draw[flowarrow] (record) -- (more3);

\draw[flowarrow]
(more3)
--
node[flowlabel]{нет}
(compare);

\draw[flowarrow]
(more3.east)
-- ++(4.7cm,0)
node[
  flowlabel,
  above,
  pos=0.18
]{да}
|- ([xshift=7mm]take.east)
-- (take.east);

\draw[flowarrow] (compare) -- (e3);

\end{tikzpicture}
\end{center}

\clearpage

% =========================================================
% ИТОГОВЫЙ ЧЕК-ЛИСТ
% =========================================================

\topic{5. Итоговый чек-лист перед решением}

\begin{enumerate}

  \item
  Смешанные числа переведены в неправильные дроби.

  \item
  При делении проверено, что делитель не равен нулю,
  и использована обратная дробь.

  \item
  Для сложения и вычитания выбран удобный общий знаменатель.

  \item
  Если использована формула
  \[
  \frac{ad\pm bc}{bd},
  \]
  итоговая дробь проверена на возможность сокращения.

  \item
  При делении десятичных дробей запятая убрана
  прежде всего из делителя умножением обоих чисел
  на одно и то же \(10^k\).

  \item
  При подстановке отрицательные значения записаны в скобках.

  \item
  На числовой прямой учтены знак,
  положение относительно нуля и модуль.

  \item
  Для сравнения дробей выбран единый удобный формат.

  \item
  Для \(\sqrt{x}\) найдены соседние полные квадраты.

  \item
  Ответ записан в требуемом формате:
  число, дробь или номер варианта.

\end{enumerate}

% =========================================================
% ТРЕНИРОВОЧНЫЙ БЛОК
% =========================================================

\clearpage
\thispagestyle{plain}

\begin{center}
{\LARGE\bfseries
\textcolor{darkaccent}{Тренировочный блок}}
\\[0.25em]
{\large Путь А --- по нарастающей сложности}
\end{center}

\vspace{0.35cm}

{\small

\textbf{Средний уровень}

\begin{enumerate}[itemsep=0.55em,topsep=0.3em]

  \item
  Вычислите:
  \[
  \frac{3}{20}+\frac{7}{25}.
  \]

  \item
  Вычислите:
  \[
  5\frac14-3\frac{7}{10}.
  \]

  \item
  Вычислите:
  \[
  \frac{9}{14}:\frac{3}{7}.
  \]

\end{enumerate}

\textbf{Выше среднего}

\begin{enumerate}[resume,itemsep=0.6em,topsep=0.3em]

  \item
  Вычислите:
  \[
  \frac{1\frac35}{\frac{4}{25}}.
  \]

  \item
  Вычислите:
  \[
  2{,}52:0{,}07.
  \]

  \item
  Вычислите удобным способом и сократите ответ:
  \[
  \frac{1}{30}+\frac{1}{42}.
  \]

  \item
  На числовой прямой точка \(a\) соответствует числу,
  приблизительно равному \(-2{,}4\).
  Какое утверждение верно?

  \[
  \begin{aligned}
  1)&\ a+2>0; &
  2)&\ -1-a>0;\\
  3)&\ 2a+5<0; &
  4)&\ -a<2.
  \end{aligned}
  \]

  \item
  Известно, что
  \[
  a\approx-3{,}2,
  \qquad
  b\approx1{,}4.
  \]
  Какое утверждение верно?

  \[
  \begin{aligned}
  1)&\ a+b>0; &
  2)&\ b-a<0;\\
  3)&\ ab<0; &
  4)&\ a^2<b^2.
  \end{aligned}
  \]

  \item
  Расположите числа в порядке возрастания:
  \[
  \frac38,
  \qquad
  \frac13,
  \qquad
  0{,}36,
  \qquad
  \frac25.
  \]

\end{enumerate}

\textbf{Сложный уровень}

\begin{enumerate}[resume,itemsep=0.5em,topsep=0.3em]

  \item
  Расположите числа в порядке возрастания:
  \[
  7{,}2,
  \qquad
  \sqrt{53},
  \qquad
  \frac{22}{3},
  \qquad
  7{,}35.
  \]

\end{enumerate}

}

% =========================================================
% ОТВЕТЫ
% =========================================================

\clearpage
\thispagestyle{plain}

\begin{center}
{\LARGE\bfseries
\textcolor{darkaccent}{Ответы к тренировочному блоку}}
\end{center}

\vspace{0.7cm}

\begin{table}[ht]
\centering
\caption{Ответы}

\begin{tabularx}{\linewidth}{
  @{}
  >{\bfseries}c
  X
  @{}
}
\toprule
№ & Ответ \\
\midrule

1 &
\(\displaystyle
\frac{43}{100}=0{,}43
\)
\\

2 &
\(\displaystyle
\frac{31}{20}=1{,}55
\)
\\

3 &
\(\displaystyle
\frac32=1{,}5
\)
\\

4 &
\(10\)
\\

5 &
\(36\)
\\

6 &
\(\displaystyle
\frac{2}{35}
\)
\\

7 &
\(2\)
\\

8 &
\(3\)
\\

9 &
\(\displaystyle
\frac13<0{,}36<\frac38<\frac25
\)
\\

10 &
\(\displaystyle
7{,}2<\sqrt{53}<\frac{22}{3}<7{,}35
\)
\\

\bottomrule
\end{tabularx}
\end{table}

\end{document}